% NeurIPS 2026 anonymous-submission paper draft.
\documentclass{article}

\PassOptionsToPackage{numbers,sort&compress}{natbib}
\PassOptionsToPackage{hypertexnames=false}{hyperref}
\usepackage{neurips_2026}

\usepackage[dvipsnames]{xcolor}
\definecolor{Bleu}{RGB}{0,0,204}
\makeatletter
\@ifpackageloaded{hyperref}{}{\usepackage{hyperref}}
\makeatother
\hypersetup{
  colorlinks=true,
  citecolor=Bleu,
  linkcolor=Bleu,
  urlcolor=Bleu,
  hypertexnames=false
}

\usepackage{amsmath,amssymb,amsthm,mathtools}
\usepackage{booktabs}
\usepackage{bbm}
\usepackage{enumitem}
\usepackage{graphicx}
\usepackage{microtype}
\usepackage{algorithm}
\usepackage{algorithmic}

\DeclareMathOperator*{\argmin}{arg\,min}
\DeclareMathOperator*{\argmax}{arg\,max}
\DeclareMathOperator{\diag}{diag}
\DeclareMathOperator{\spann}{span}
\DeclareMathOperator{\rank}{rank}
\newcommand{\E}{\mathbb{E}}
\newcommand{\R}{\mathbb{R}}
\newcommand{\cS}{\mathcal S}
\newcommand{\cA}{\mathcal A}
\newcommand{\cX}{\mathcal X}
\newcommand{\cD}{\mathcal D}
\newcommand{\cT}{\mathcal T}
\newcommand{\cF}{\mathcal F}
\newcommand{\cV}{\mathcal V}
\newcommand{\one}{\mathbbm 1}
\newcommand{\norm}[1]{\left\lVert #1 \right\rVert}
\newcommand{\linf}[1]{\left\lVert #1 \right\rVert_\infty}
\newcommand{\hatP}{\widehat P}
\newcommand{\hatr}{\widehat r}
\newcommand{\hatq}{\widehat q}
\newcommand{\hatQ}{\widehat Q}
\newcommand{\hatA}{\widehat A}
\newcommand{\hatb}{\widehat b}
\newcommand{\tildeA}{\widetilde A}
\newcommand{\tildeb}{\widetilde b}
\newcommand{\bell}{\mathcal B}
\newcommand{\proj}{\Pi}

\theoremstyle{plain}
\newtheorem{theorem}{Theorem}
\newtheorem{lemma}{Lemma}
\newtheorem{corollary}{Corollary}
\newtheorem{proposition}{Proposition}
\theoremstyle{definition}
\newtheorem{definition}{Definition}
\newtheorem{assumption}{Assumption}
\theoremstyle{remark}
\newtheorem{remark}{Remark}

\title{Target-Weighted Bellman Aggregation Trees and Forests for Offline Policy Evaluation}

\author{%
  Anonymous Authors
}

\begin{document}

\maketitle

\begin{abstract}
Regression trees and random forests are natural tools for interpretable and
adaptive value estimation, but tree-based fitted value methods usually treat
Bellman updates as ordinary supervised labels. The resulting splits are
optimized for transient pseudo-outcomes and for the sampling distribution of
the offline data, which need not match the distribution under which value
accuracy matters. We propose \emph{target-weighted Bellman Aggregation Trees
and Forests} for offline policy evaluation. A tree defines a partition, while a
forest defines a finite collection of leaf-indicator features. In both cases,
the final value estimate is obtained by solving one target-weighted projected
Bellman equation, rather than by iterating supervised regressions. Single
Bellman Aggregation Trees (BATs) give the interpretable partition special case;
Bellman Aggregation Forests (BAFs) use randomized honest trees to build a
richer finite feature space before a joint Bellman solve. For fixed or honest
learned features, the population estimator is the projected Bellman fixed
point, with approximation error controlled by the projection error of
\(Q^\pi\). Finite samples, rank reduction, validation pruning, and estimated
weights enter as perturbations of the projected Bellman linear system. The
manuscript gives a method and theory draft with a falsifiable experimental
protocol; empirical BAT/BAF results remain pending.
\end{abstract}

\section{Introduction}

Tree-based value estimation is appealing for a simple reason: a tree gives a
readable partition of the state-action space. If a fitted value function
assigns the same number to all points in a leaf, one would like that leaf to
mean something dynamical: points in the same leaf should have similar immediate
rewards and similar continuation behavior under the policy being evaluated, in
the distribution where accuracy matters. Standard tree-based fitted value
methods do not enforce this interpretation. They usually embed regression
trees or tree ensembles inside fitted Bellman iteration: at each round, a
Bellman target is computed from the previous value estimate and a supervised
regressor is fit to that target \citep{ernst2005tree,munos2008finite}. This is
a useful modular algorithm, but the fitted tree is optimized for a moving
pseudo-label, usually under the behavior distribution. A split may reduce
prediction error for the current target while merging states whose
target-policy transitions lead to different parts of the value function.

This paper starts from the opposite perspective. A tree is not only a
supervised learner; it is also a finite aggregation of the state-action space.
A random forest is not only an average of supervised trees; it is also a
finite, adaptive feature map built from many leaf indicators. Once such a
feature map and an evaluation distribution \(\mu\) are fixed, policy evaluation
has a canonical solution: project the Bellman equation onto the span of those
features under \(\mu\). For a single tree this projection is exactly the
aggregate Bellman equation between leaves. For a forest it is a single linear
projected Bellman solve over all forest leaf features.

We propose \emph{target-weighted Bellman Aggregation Trees and Forests}. Given
offline transitions \((S_i,A_i,R_i,S'_i)\) from a behavior distribution
\(\nu\), the method learns tree or forest features over \(X=(S,A)\), evaluates
them under a reference distribution \(\mu\), typically the target policy's
stationary state-action distribution, and uses weights \(w=d\mu/d\nu\) to fit
the projected Bellman system. Splits are chosen by weighted held-out Bellman
residuals computed after provisional Bellman solves. The final value estimate
is then refit on an independent estimation sample. For BAF, all provisional
per-tree values are discarded; the retained object is the forest feature map,
and the final values come from one joint projected Bellman solve.

The contribution is deliberately narrow. BATs and BAFs are not proposed as deep
offline RL systems. They address a more basic statistical question: how should
one fit regression trees and random forests when the regression target is a
Bellman fixed point and the evaluation distribution differs from the sampling
distribution? The answer is to treat tree leaves as Bellman aggregation
features, then solve the target-weighted projected Bellman equation induced by
those features.

Our contributions are:
\begin{itemize}[leftmargin=1.5em]
\item We formulate BATs and BAFs as non-iterative estimators for offline
policy evaluation: a single tree is a partition-indicator feature map, and a
forest is a randomized collection of leaf-indicator features.
\item We give a practical construction using weighted Bellman-aware split
scores, honest sample splitting, stabilized weights, validation pruning, and
honest rank reduction before the final forest solve.
\item We prove that, for fixed or honest learned features, the estimator solves
a projected Bellman equation and has aggregation bias controlled by
\(\|(I-\proj^\mu_\psi)Q^\pi\|_{L^2(\mu)}\). For trees this is within-leaf
variation; for forests it is projection error onto the forest leaf-feature
span.
\item We state finite-sample and weight-estimation perturbation bounds for the
projected Bellman linear system, plus a finite-candidate validation guarantee.
\item We give an experimental protocol, with results pending, designed to
isolate whether Bellman-aware tree and forest features improve value accuracy
per degree of freedom compared with ordinary supervised tree objectives.
\end{itemize}

\section{Problem Setup}
\label{sec:setup}

We consider an infinite-horizon discounted Markov decision process with state
space \(\cS\), finite action space \(\cA\), transition kernel \(K\), bounded
reward \(|R|\le R_{\max}\), and discount \(\gamma\in(0,1)\). The target policy
\(\pi(a\mid s)\) is fixed. Let \(X=(S,A)\in\cX=\cS\times\cA\), and define
\[
  X^\pi_+ = (S',A'),\qquad
  S'\sim K(\cdot\mid S,A),\quad A'\sim \pi(\cdot\mid S').
\]
The policy-evaluation Bellman operator is
\[
  (T^\pi Q)(s,a)
  =
  \E\!\left[R+\gamma Q(S',A')\mid S=s,A=a,\; A'\sim\pi(\cdot\mid S')\right],
\]
and \(Q^\pi\) is its unique fixed point. We observe \(n\) i.i.d. one-step
offline transitions \((S_i,A_i,R_i,S'_i)\) with
\(X_i=(S_i,A_i)\sim\nu\). Trajectory data require standard mixing or blocking
arguments. We evaluate error under a reference distribution \(\mu\) on
\(\cX\), typically the target-policy stationary state-action distribution
\(\mu_\pi\), and assume \(\mu\ll\nu\) with density ratio \(w=d\mu/d\nu\).

\paragraph{Finite Bellman features.}
Let \(\psi:\cX\to\R^d\) be a finite feature map selected independently of the
final estimation sample. It represents value functions
\[
  Q_\theta(x)=\psi(x)^\top\theta,\qquad \theta\in\R^d.
\]
For a next state \(s'\), write
\[
  \bar\psi(s')=\E_{A'\sim\pi(\cdot\mid s')}[\psi(s',A')].
\]
The target-weighted projected Bellman equations are
\[
  A^\mu_\psi\theta=b^\mu_\psi,\qquad
  A^\mu_\psi
  =
  \E_\mu\!\left[\psi(X)\{\psi(X)-\gamma\psi(X^\pi_+)\}^\top\right],
  \quad
  b^\mu_\psi
  =
  \E_\mu[\psi(X)R].
  \label{eq:population_linear_system}
\]
With behavior data, these moments are estimated by importance weighting:
\[
  \hatA_\psi
  =
  \frac{\sum_i \hat w_i\psi_i(\psi_i-\gamma\bar\psi_i^+)^\top}
       {\sum_i\hat w_i},
  \qquad
  \hatb_\psi
  =
  \frac{\sum_i\hat w_i\psi_i R_i}{\sum_i\hat w_i},
  \label{eq:empirical_linear_system}
\]
where \(\psi_i=\psi(S_i,A_i)\) and
\(\bar\psi_i^+=\bar\psi(S'_i)\). The final estimate is
\(\hat Q_\psi(x)=\psi(x)^\top\hat\theta\), where
\(\hatA_\psi\hat\theta=\hatb_\psi\) after deterministic rank reduction if
needed.

\paragraph{Trees and forests as features.}
A regression tree \(\cT\) partitions \(\cX\) into leaves
\(\ell\in\{1,\ldots,L\}\). Its feature map is
\(\psi_{\cT}(x)=e_{\phi_{\cT}(x)}\), the one-hot indicator of the leaf
containing \(x\). In this case \eqref{eq:population_linear_system} is exactly
the aggregate Bellman equation
\[
  q=r^\mu+\gamma P^\mu q,
  \qquad
  r^\mu_\ell=\E_\mu[R\mid X\in\ell],
  \quad
  P^\mu_{\ell m}=\Pr_\mu\{\phi_{\cT}(X^\pi_+)=m\mid X\in\ell\}.
\]
A random forest \(\cF=(\cT_1,\ldots,\cT_B)\) defines the concatenated leaf
feature map
\[
  \psi_{\cF}(x)
  =
  \frac{1}{B}\big(e_{\phi_{\cT_1}(x)}^\top,\ldots,
  e_{\phi_{\cT_B}(x)}^\top\big)^\top .
\]
The scaling is fixed and immaterial after rank reduction; it keeps the forest
features on a comparable numerical scale as \(B\) changes. BAF solves
\eqref{eq:empirical_linear_system} once over the span of these forest
features. It is not the average of \(B\) independently labeled BATs.

\section{Bellman Aggregation Trees and Forests}
\label{sec:method}

The method has two stages. First, it constructs tree or forest features using
Bellman-aware split scores. Second, it discards provisional split-time values
and solves the final target-weighted projected Bellman system on an independent
estimation sample. The theorem-clean version uses five disjoint sample roles:
growth-fit, growth-score, validation, rank, and estimation. Cross-fitting can
recover data efficiency, but honesty keeps the estimand transparent.

\subsection{Bellman-Aware Split Scoring}

For a candidate feature map \(\psi\), estimate \(\hatA_\psi,\hatb_\psi\) on
the growth-fit sample and solve
\(\hatA_\psi\hat\theta_\psi=\hatb_\psi\). Score the candidate on a held-out
sample by
\[
  \delta_i(\psi)
  =
  R_i+\gamma(\bar\psi_i^+)^\top\hat\theta_\psi
  -\psi_i^\top\hat\theta_\psi,
  \qquad
  \mathrm{Score}_\mu(\psi)
  =
  \frac{\sum_{i\in\cD_{\mathrm{hold}}}\hat w_i\delta_i(\psi)^2}
       {\sum_{i\in\cD_{\mathrm{hold}}}\hat w_i}
  +\lambda\,\dim(\spann\psi).
  \label{eq:score}
\]
During growth, \(\cD_{\mathrm{hold}}\) is the growth-score sample. The
validation sample is used only to prune a finite path or finite forest
candidate set.

For a single BAT, the candidate maps are one-hot leaf indicators for candidate
tree splits, and the provisional solve can be implemented as the aggregate
leaf Bellman solve. For BAF, randomized BAT-style trees are grown using
subsampling and feature subsampling. The split score is a feature-construction
criterion; the final forest values are obtained only after all trees have been
grown and the joint forest feature map has been rank reduced.

\begin{algorithm}[t]
\caption{Target-Weighted Bellman Aggregation Forest}
\label{alg:baf}
\begin{algorithmic}[1]
\REQUIRE Offline data \(\cD\), target policy \(\pi\), raw weights
\(\tilde w\), number of trees \(B\), discount \(\gamma\), split and pruning
parameters, stabilization parameters.
\STATE Split \(\cD\) into growth-fit, growth-score, validation, rank, and estimation samples.
\STATE Stabilize and normalize \(\tilde w\) within each sample to obtain \(\hat w\).
\FOR{\(b=1,\ldots,B\)}
  \STATE Draw randomized growth-fit and growth-score subsamples and feature subsets.
  \STATE Grow a BAT-style tree path using admissible splits with minimum raw count, weighted mass, and ESS.
  \STATE Score candidate splits by \eqref{eq:score} using provisional Bellman solves on the current tree.
  \STATE Prune the finite path on the validation sample and store the selected tree \(\cT_b\).
\ENDFOR
\STATE Build \(\psi_{\cF}\) from all selected forest leaves; for BAT take \(B=1\).
\STATE Drop linearly dependent columns by deterministic QR/SVD on the rank sample.
\STATE Estimate \(\hatA_{\psi_{\cF}}\) and \(\hatb_{\psi_{\cF}}\) by \eqref{eq:empirical_linear_system}.
\STATE Solve \(\hatA_{\psi_{\cF}}\hat\theta=\hatb_{\psi_{\cF}}\).
\STATE \textbf{return} \(\hat Q(x)=\psi_{\cF}(x)^\top\hat\theta\).
\end{algorithmic}
\end{algorithm}

\subsection{Stabilized Weights and Rank Reduction}

Given raw density-ratio estimates \(\tilde w_i\), the practical estimator clips
and shrinks them before tree construction:
\[
  \bar w_i=(1-\alpha)\min\{\tilde w_i,c\}+\alpha,\qquad
  \hat w_i=\frac{\bar w_i}{|\cD_s|^{-1}\sum_{j\in\cD_s}\bar w_j},
  \quad i\in\cD_s .
\]
Here \(\cD_s\) is one of the sample roles. A leaf is admissible only if it
has enough raw observations, enough weighted mass, and effective sample size
\[
  \mathrm{ESS}(\ell)
  =
  \frac{\left(\sum_{i:\phi(X_i)=\ell}\hat w_i\right)^2}
       {\sum_{i:\phi(X_i)=\ell}\hat w_i^2}
\]
above a prespecified threshold. Forest features are often linearly dependent:
each tree contributes leaf indicators that sum to a constant, and different
trees may create redundant columns. The theorem-clean implementation performs
deterministic rank reduction on an independent rank sample before the projected
Bellman solve. Using the estimation design for QR/SVD is a pragmatic
implementation shortcut because it uses no rewards, but the stated
finite-sample theorem treats the retained full-rank feature map as fixed before
final estimation. Ridge regularization is reported separately because it
changes the estimand.

\section{Theory}
\label{sec:theory}

The theory treats a learned tree or forest as fixed after conditioning on the
growth, validation, and rank-reduction data. The central object is the feature span
\(\cV_\psi=\{\psi^\top\theta:\theta\in\R^d\}\), not the particular algorithm
that produced it.
After rank reduction, \(d=d_r=\dim(\cV_\psi)\) denotes the retained feature
dimension. We write
\[
  (P_\pi f)(x)=\E[f(X^\pi_+)\mid X=x],
  \qquad
  \proj^\mu_\psi f
  =
  \argmin_{g\in\cV_\psi}\|f-g\|_{L^2(\mu)}
\]
for the target-policy transition operator and the \(L^2(\mu)\) orthogonal
projection onto the retained feature span.

\begin{assumption}[Sampling, overlap, and boundedness]
\label{ass:coverage}
The final estimation sample consists of i.i.d. one-step transitions with
\(X_i\sim\nu\). The reference measure satisfies \(\mu\ll\nu\), the oracle
weight \(w=d\mu/d\nu\) is known for the main finite-sample theorem and bounded
by \(W\), rewards satisfy \(|R|\le R_{\max}\), and \(\gamma\in(0,1)\).
\end{assumption}

\begin{assumption}[Reference-measure contraction]
\label{ass:l2contract}
The Bellman operator is a \(\gamma\)-contraction in \(L^2(\mu)\):
\[
  \|T^\pi f-T^\pi g\|_{L^2(\mu)}
  \le
  \gamma\|f-g\|_{L^2(\mu)}
  \qquad\text{for all bounded } f,g .
\]
This holds, for example, when \(\mu=\mu_\pi\) is a stationary state-action
distribution for the target-policy Markov chain.
\end{assumption}

\begin{lemma}[Stationary reference measures are contraction measures]
\label{lem:stationary_contraction}
If \(\mu_\pi\) is stationary for the target-policy Markov chain on
\(\cX\), then \(T^\pi\) is a \(\gamma\)-contraction in \(L^2(\mu_\pi)\).
\end{lemma}

\begin{assumption}[Honest finite features and rank]
\label{ass:rank}
The feature map \(\psi:\cX\to\R^d\) is selected independently of the final
estimation sample, satisfies \(\|\psi(x)\|_2\le K\), and has full rank after
deterministic rank reduction. The population projected Bellman matrix
\(A^\mu_\psi\) in \eqref{eq:population_linear_system} has
\(\sigma_{\min}(A^\mu_\psi)\ge \kappa>0\).
\end{assumption}

\begin{proposition}[Projected Bellman target]
\label{prop:projected}
Under Assumption~\ref{ass:l2contract}, the operator
\(\proj^\mu_\psi T^\pi\), restricted to \(\cV_\psi\), is a
\(\gamma\)-contraction, where \(\proj^\mu_\psi\) is the \(L^2(\mu)\) projection
onto \(\cV_\psi\). Hence there is a unique \(Q^\mu_\psi\in\cV_\psi\) satisfying
\[
  Q^\mu_\psi=\proj^\mu_\psi T^\pi Q^\mu_\psi .
\]
If \(A^\mu_\psi\) is nonsingular, then
\(Q^\mu_\psi(x)=\psi(x)^\top\theta^\mu_\psi\) with
\(A^\mu_\psi\theta^\mu_\psi=b^\mu_\psi\).
\end{proposition}

\begin{theorem}[Aggregation bias for trees and forests]
\label{thm:aggregation_bias}
Under Assumption~\ref{ass:l2contract}, any fixed tree or forest feature map
\(\psi\) satisfies
\[
  \|Q^\mu_\psi-Q^\pi\|_{L^2(\mu)}
  \le
  \frac{1}{1-\gamma}
  \|(I-\proj^\mu_\psi)Q^\pi\|_{L^2(\mu)}.
\]
For a BAT, let \(\proj^\mu_{\cT}\) denote the \(L^2(\mu)\) projection onto the
tree leaves. The right-hand side is the within-leaf variation of \(Q^\pi\)
under \(\mu\):
\[
  \|(I-\proj^\mu_{\cT})Q^\pi\|_{L^2(\mu)}^2
  =
  \E_\mu[\operatorname{Var}\{Q^\pi(X)\mid \phi_{\cT}(X)\}].
\]
For a BAF, it is the approximation error of \(Q^\pi\) by the span of all
retained forest leaf features.
\end{theorem}

\begin{proposition}[Forest feature monotonicity]
\label{prop:monotonicity}
If \(\cV_{\psi_1}\subseteq \cV_{\psi_2}\), then
\[
  \|(I-\proj^\mu_{\psi_2})Q^\pi\|_{L^2(\mu)}
  \le
  \|(I-\proj^\mu_{\psi_1})Q^\pi\|_{L^2(\mu)}.
\]
Thus, after rank reduction, adding linearly independent forest leaf features
cannot increase the population approximation term, although the finite-sample
error depends on the resulting dimension and conditioning of \(A^\mu_\psi\).
\end{proposition}

\begin{theorem}[Projected Bellman perturbation]
\label{thm:feature_perturbation}
Let \(\theta=\theta^\mu_\psi\) and \(\hat\theta\) solve
\(A^\mu_\psi\theta=b^\mu_\psi\) and
\(\hatA_\psi\hat\theta=\hatb_\psi\). If
\[
  \|\hatA_\psi-A^\mu_\psi\|_{\mathrm{op}}\le \varepsilon_A<\kappa,
  \qquad
  \|\hatb_\psi-b^\mu_\psi\|_2\le \varepsilon_b,
\]
then
\[
  \|\hat\theta-\theta\|_2
  \le
  \frac{\varepsilon_b+\varepsilon_A\|\theta\|_2}{\kappa-\varepsilon_A},
  \qquad
  \|\hat Q_\psi-Q^\mu_\psi\|_{L^2(\mu)}
  \le
  K\frac{\varepsilon_b+\varepsilon_A\|\theta\|_2}{\kappa-\varepsilon_A}.
\]
\end{theorem}

\begin{proposition}[A simple concentration consequence]
\label{prop:concentration}
Assume Assumptions~\ref{ass:coverage} and~\ref{ass:rank}, with oracle weights
on the final estimation sample. Let
\[
  D_n=\frac1n\sum_{i=1}^n w_i,
  \qquad
  \tildeA_\psi=\frac1n\sum_{i=1}^n
  w_i\psi_i(\psi_i-\gamma\bar\psi_i^+)^\top,
  \qquad
  \tildeb_\psi=\frac1n\sum_{i=1}^n w_i\psi_iR_i,
\]
and \(\hatA_\psi=\tildeA_\psi/D_n\),
\(\hatb_\psi=\tildeb_\psi/D_n\). There is a universal constant \(C\) such
that, if \(n\ge 2W^2\log(6/\delta)\), then with probability at least
\(1-\delta\),
\begin{align*}
  \|\hatA_\psi-A^\mu_\psi\|_{\mathrm{op}}
  \le
  C WK^2(1+\gamma)
  \left\{
    \sqrt{\frac{d+\log(6/\delta)}{n}}
    +
    \frac{\log(6d/\delta)}{n}
  \right\},\\
  \|\hatb_\psi-b^\mu_\psi\|_2
  \le
  C WKR_{\max}
  \sqrt{\frac{d+\log(6/\delta)}{n}} .
\end{align*}
This elementary bound is not intended to be sharp; it makes explicit the
forest tradeoff between feature dimension, conditioning, and sample size.
\end{proposition}

\begin{corollary}[Tree aggregate special case]
\label{cor:tree}
For one-hot tree features, \(A^\mu_\psi=D_\mu(I-\gamma P^\mu)\) and
\(b^\mu_\psi=D_\mu r^\mu\), where
\(D_\mu=\diag(\mu_1,\ldots,\mu_L)\),
\(\mu_\ell=\mu\{\phi_{\cT}(X)=\ell\}\), and
\(\mu_{\min}=\min_\ell\mu_\ell>0\). Let
\[
  \hat\mu_\ell=\frac1n\sum_{i=1}^n
  w_i\one\{\phi_{\cT}(X_i)=\ell\}.
\]
If \(\min_\ell\hat\mu_\ell\ge \mu_{\min}/2\), and
\(\linf{\hatr^\mu-r^\mu}\le\varepsilon_r\),
\(\|\hatP^\mu-P^\mu\|_\infty\le\varepsilon_P\), then
\[
  \linf{\hatq^\mu_{\cT}-q^\mu_{\cT}}
  \le
  \frac{\varepsilon_r}{1-\gamma}
  +
  \frac{\gamma R_{\max}\varepsilon_P}{(1-\gamma)^2}.
\]
The leaf-mass condition holds with probability at least \(1-\delta\) when
\[
  n\ge 2W^2\mu_{\min}^{-2}\log(2L/\delta).
\]
\end{corollary}

\begin{proposition}[Estimated-weight perturbation]
\label{prop:weight_perturbation}
Let \(\theta^w\) and \(\theta^{\hat w}\) be the projected Bellman solutions
obtained from two weighted moment systems \((A^w,b^w)\) and
\((A^{\hat w},b^{\hat w})\). If
\(\sigma_{\min}(A^w)\ge\kappa\),
\(\|A^{\hat w}-A^w\|_{\mathrm{op}}\le\eta_A<\kappa\), and
\(\|b^{\hat w}-b^w\|_2\le\eta_b\), then
\[
  \|\theta^{\hat w}-\theta^w\|_2
  \le
  \frac{\eta_b+\eta_A\|\theta^w\|_2}{\kappa-\eta_A}.
\]
Thus estimated or stabilized weights affect BAT/BAF only through the aggregate
projected Bellman moments they induce.
\end{proposition}

\begin{proposition}[Validation pruning over a finite candidate set]
\label{prop:validation}
Let \(\mathfrak P\) be a finite set of candidate trees, forests, or feature
maps selected independently of a validation sample of size \(n_{\mathrm{val}}\).
For each \(\psi\in\mathfrak P\), let \(\delta_\psi\) be its validation
residual after the provisional Bellman solve. Suppose
\(|\delta_\psi|\le B_\delta\) and \(0\le w\le W\). If
\(\hat\psi\) minimizes the empirical weighted residual plus
\(\lambda\dim(\spann\psi)\) over \(\mathfrak P\), then with probability at
least \(1-\delta\),
\[
  L_\mu(\hat\psi)+\lambda\dim(\spann\hat\psi)
  \le
  \min_{\psi\in\mathfrak P}
  \{L_\mu(\psi)+\lambda\dim(\spann\psi)\}
  +
  2WB_\delta^2
  \sqrt{\frac{\log(2|\mathfrak P|/\delta)}{2n_{\mathrm{val}}}},
\]
where \(L_\mu(\psi)=\E_\mu[\delta_\psi^2]\).
Here the empirical risk is the unnormalized oracle-weight risk
\[
  \widehat L_{\mathrm{val}}(\psi)
  =
  \frac{1}{n_{\mathrm{val}}}
  \sum_{i\in\cD_{\mathrm{val}}}w_i\delta_{\psi,i}^2 .
\]
The self-normalized validation score obeys the same comparison up to the
standard denominator event for
\(n_{\mathrm{val}}^{-1}\sum_{i\in\cD_{\mathrm{val}}}w_i\).
\end{proposition}

Combining the population and sample terms gives
\[
  \|\hat Q_\psi-Q^\pi\|_{L^2(\mu)}
  \le
  \underbrace{\|Q^\mu_\psi-Q^\pi\|_{L^2(\mu)}}_{\textnormal{projection bias}}
  +
  \underbrace{\|\hat Q_\psi-Q^\mu_\psi\|_{L^2(\mu)}}_{\textnormal{statistical error}}.
\]
Honesty lets this bound be read conditionally on the learned tree or forest.

\section{Experimental Protocol}
\label{sec:experiments}

The empirical part of this manuscript remains a protocol rather than a
completed results section. The protocol tests the mechanism: whether
target-weighted Bellman-aware tree and forest features improve value accuracy
per degree of freedom when supervised prediction boundaries and Bellman-relevant
transition boundaries differ.

\paragraph{Diagnostic MDP.}
Use a small continuous-state MDP where immediate reward variation and
continuation-value variation occur along different axes. A supervised tree or
forest fit to Monte Carlo returns or one-step pseudo-labels can split on reward
while merging states with different target-policy continuations. The diagnostic
figure reports the true \(Q^\pi\), learned BAT partitions, forest-induced
nearest-neighbor/feature weights, and scalar errors.

\paragraph{Grid-auditable navigation.}
Use a two-dimensional navigation task with discrete actions, nonlinear drift,
and offline data whose behavior distribution under-samples regions important
under \(\mu_\pi\). A fine grid solver supplies reference values and stationary
norm diagnostics; the learning algorithms see only offline transitions.

\paragraph{Baselines and ablations.}
Compare BAT and BAF against CART on Monte Carlo returns, random-forest returns,
Tree-FQE, RF-FQE, random tree partitions, and uniform grids. The main ablation
compares supervised split scores with Bellman-aware split scores. A second
ablation compares the joint BAF solve with bagged independently solved BATs,
which tests whether the projected forest formulation matters.

\paragraph{Metrics.}
Report stationary \(L^2(\mu_\pi)\) value error, weighted Bellman residual,
policy-value error, error versus number of leaves or retained features, and
error versus number of trees. The target pattern is lower value error per
degree of freedom and partitions or forest weights aligned with
transition-induced value boundaries. Residuals are diagnostic, not a substitute
for value error under \(\mu_\pi\).

\section{Related Work}
\label{sec:related}

\paragraph{Tree-based fitted value methods.}
Tree-based batch-mode reinforcement learning uses regression trees and tree
ensembles inside fitted value iteration \citep{ernst2005tree}. The
policy-evaluation analogue, Tree-FQE or RF-FQE, treats trees as supervised
regressors for Bellman pseudo-labels. BATs and BAFs instead use trees to build
aggregation features and compute values by a projected Bellman solve.

\paragraph{Random forests and forest features.}
Random forests average randomized trees to obtain adaptive nonparametric
regressors \citep{breiman2001random}. Modern work also views forests through
adaptive kernels, feature spans, and honesty
\citep{wager2018estimation,athey2019generalized,scornet2016random}. BAF uses
the forest as an adaptive Bellman feature map: the final estimate is not a
supervised forest average, but the projected fixed point over all retained leaf
features.

\paragraph{State aggregation and projected Bellman equations.}
State aggregation reduces dynamic programming to a finite abstract process
\citep{van2006performance,bertsekas2011approximate}. Least-squares temporal
difference methods solve projected Bellman equations under linear function
approximation \citep{tsitsiklis1996analysis,lagoudakis2003least}. BAT is the
tree-partition instance of this idea; BAF is the random-forest feature instance.

\paragraph{Bellman residual and minimax methods.}
Bellman residual minimization has a long history \citep{baird1995residual}.
Modern offline RL methods use minimax or adversarial Bellman-error objectives
to avoid strong completeness assumptions
\citep{xie2020q,uehara2021finite,xie2021bellman}. BAT/BAF use residuals in a
restricted way: residuals construct and prune features, while final values come
from the projected Bellman equation.

\section{Discussion and Limitations}
\label{sec:discussion}

BATs and BAFs answer a narrow question: how should tree and forest regressors
be fit when the target is a Bellman fixed point and the evaluation distribution
differs from the sampling distribution? The proposed answer is to use trees as
target-weighted aggregation features and solve the induced projected Bellman
equation. Single trees retain an interpretable partition; forests trade some
of that direct interpretability for a richer approximation space.

The paper focuses on fixed-policy evaluation. For optimal control, the Bellman
equation contains a max over actions, so the final linear solve no longer
applies. A second limitation is conditioning: rich forest feature spaces can
reduce approximation bias while increasing variance and numerical
ill-conditioning, which is why deterministic rank reduction is part of the
estimator. Finally, this manuscript is evidence-incomplete. It specifies the
experiments needed to test the mechanism, but does not yet report BAT/BAF
results.

\section*{Reproducibility Statement}

The implementation uses deterministic data splits, including a rank-reduction
split for the theorem-clean version, recorded random seeds for forest
subsampling and feature subsampling, declared weight stabilization, minimum
raw-count, weighted-mass and ESS thresholds, validation-based pruning, and
deterministic rank reduction before the final solve. Experiments report the
full tree-size and forest-size paths, not only the selected model.

\bibliographystyle{plainnat}
\bibliography{../../common/ref}

\appendix

\section{Proofs}
\label{app:proofs}

\paragraph{Standard facts used below.}
We use the Hilbert-space projection theorem and Banach fixed-point theorem
\citep{conway1994course}, Weyl's singular-value perturbation inequality
\citep{horn2012matrix}, Hoeffding's inequality for bounded random variables
\citep{hoeffding1963probability}, matrix Bernstein
\citep{tropp2012user}, and standard vector concentration for bounded random
vectors \citep{wainwright2019high}.

\begin{proof}[Proof of Lemma~\ref{lem:stationary_contraction}]
For any bounded \(f,g\), rewards cancel and
\[
  T^\pi f(X)-T^\pi g(X)
  =
  \gamma (P_\pi(f-g))(X).
\]
By Jensen's inequality,
\[
  \{P_\pi(f-g)(X)\}^2
  \le
  P_\pi\{(f-g)^2\}(X).
\]
Taking \(X\sim\mu_\pi\) and using stationarity of \(\mu_\pi\) for the
target-policy chain gives
\[
  \|P_\pi(f-g)\|_{L^2(\mu_\pi)}^2
  \le
  \E_{\mu_\pi}[(f-g)^2(X^\pi_+)]
  =
  \|f-g\|_{L^2(\mu_\pi)}^2.
\]
Multiplying by \(\gamma^2\) proves the claim.
\end{proof}

\begin{proof}[Proof of Proposition~\ref{prop:projected}]
Because \(\cV_\psi\) is finite dimensional, it is a closed subspace of
\(L^2(\mu)\). The orthogonal projection
\(\proj^\mu_\psi:L^2(\mu)\to\cV_\psi\) is therefore well defined,
maps every function into \(\cV_\psi\), and is nonexpansive. For
\(f,g\in\cV_\psi\),
\[
  \|\proj^\mu_\psi T^\pi f-\proj^\mu_\psi T^\pi g\|_{L^2(\mu)}
  \le
  \|T^\pi f-T^\pi g\|_{L^2(\mu)}
  \le
  \gamma\|f-g\|_{L^2(\mu)}.
\]
Thus \(\proj^\mu_\psi T^\pi\) is a contraction on the complete metric space
\(\cV_\psi\), so Banach's theorem gives a unique fixed function
\(Q^\mu_\psi\in\cV_\psi\).

If \(Q(x)=\psi(x)^\top\theta\), the projected fixed-point equation is
equivalent to the residual \(Q-T^\pi Q\) being orthogonal to every element of
\(\cV_\psi\). Since the coordinates of \(\psi\) span \(\cV_\psi\), this is
equivalent to
\[
  \E_\mu[\psi(X)\{Q(X)-R-\gamma Q(X^\pi_+)\}]=0,
\]
which is exactly \(A^\mu_\psi\theta=b^\mu_\psi\). The fixed function is unique
by contraction. The coefficient vector is unique whenever
\(A^\mu_\psi\) is nonsingular, which is enforced after rank reduction in
Assumption~\ref{ass:rank}.
\end{proof}

\begin{proof}[Proof of Theorem~\ref{thm:aggregation_bias}]
Using Proposition~\ref{prop:projected}, \(Q^\pi=T^\pi Q^\pi\), nonexpansiveness
of \(\proj^\mu_\psi\), and Assumption~\ref{ass:l2contract},
\[
  \begin{aligned}
  \|Q^\mu_\psi-Q^\pi\|_{L^2(\mu)}
  &=
  \|\proj^\mu_\psi T^\pi Q^\mu_\psi-T^\pi Q^\pi\|_{L^2(\mu)}\\
  &\le
  \|\proj^\mu_\psi(T^\pi Q^\mu_\psi-T^\pi Q^\pi)\|_{L^2(\mu)}
  +
  \|(I-\proj^\mu_\psi)Q^\pi\|_{L^2(\mu)}\\
  &\le
  \gamma\|Q^\mu_\psi-Q^\pi\|_{L^2(\mu)}
  +
  \|(I-\proj^\mu_\psi)Q^\pi\|_{L^2(\mu)}.
  \end{aligned}
\]
Rearranging proves the bound.
For one-hot tree features, \(\proj^\mu_{\cT}Q^\pi\) is the conditional mean
\(\E_\mu[Q^\pi(X)\mid \phi_{\cT}(X)]\). The displayed variance identity is
therefore the conditional-variance decomposition
\[
  \E_\mu\!\left[
    \{Q^\pi(X)-\E_\mu[Q^\pi(X)\mid \phi_{\cT}(X)]\}^2
  \right]
  =
  \E_\mu[\operatorname{Var}\{Q^\pi(X)\mid \phi_{\cT}(X)\}].
\]
\end{proof}

\begin{proof}[Proof of Proposition~\ref{prop:monotonicity}]
Orthogonal projection onto a larger closed subspace cannot have larger
least-squares approximation error. Since \(\cV_{\psi_1}\subseteq\cV_{\psi_2}\),
the best approximation to \(Q^\pi\) in \(\cV_{\psi_2}\) is at least as good as
the best approximation in \(\cV_{\psi_1}\).
\end{proof}

\begin{proof}[Proof of Theorem~\ref{thm:feature_perturbation}]
Write \(A=A^\mu_\psi\) and \(b=b^\mu_\psi\). Subtract
\(A\theta=b\) from \(\hatA\hat\theta=\hatb\):
\[
  \hatA(\hat\theta-\theta)
  =
  \hatb-b-(\hatA-A)\theta .
\]
Weyl's inequality gives
\(\sigma_{\min}(\hatA)\ge\kappa-\varepsilon_A\), hence
\[
  \|\hat\theta-\theta\|_2
  \le
  \frac{\varepsilon_b+\varepsilon_A\|\theta\|_2}{\kappa-\varepsilon_A}.
\]
The value-function bound follows from
\(\|\psi^\top(\hat\theta-\theta)\|_{L^2(\mu)}
\le \sup_x\|\psi(x)\|_2\|\hat\theta-\theta\|_2\).
\end{proof}

\begin{proof}[Proof of Proposition~\ref{prop:concentration}]
Condition on the fixed retained feature map. Define
\[
  Y_i^A=w_i\psi_i(\psi_i-\gamma\bar\psi_i^+)^\top,
  \qquad
  Y_i^b=w_i\psi_iR_i.
\]
Then \(\E_\nu Y_i^A=A^\mu_\psi\) and
\(\E_\nu Y_i^b=b^\mu_\psi\). Moreover
\[
  \|Y_i^A\|_{\mathrm{op}}\le WK^2(1+\gamma),
  \qquad
  \|Y_i^b\|_2\le WKR_{\max}.
\]
Matrix Bernstein and vector Bernstein imply that, with probability at least
\(1-2\delta/3\),
\[
  \|\tildeA_\psi-A^\mu_\psi\|_{\mathrm{op}}
  \le
  C_0 WK^2(1+\gamma)
  \left\{
    \sqrt{\frac{d+\log(6/\delta)}{n}}
    +
    \frac{\log(6d/\delta)}{n}
  \right\},
\]
and
\[
  \|\tildeb_\psi-b^\mu_\psi\|_2
  \le
  C_0 WKR_{\max}
  \sqrt{\frac{d+\log(6/\delta)}{n}} .
\]
Since \(0\le w_i\le W\) and \(\E_\nu w_i=1\), Hoeffding's inequality gives
\[
  \Pr\{|D_n-1|>t\}
  \le
  2\exp\{-2nt^2/W^2\}.
\]
With \(t_n=W\sqrt{\log(6/\delta)/(2n)}\), this event has probability at
least \(1-\delta/3\). In particular,
\[
  \Pr\{|D_n-1|>1/2\}
  \le
  2\exp\{-n/(2W^2)\},
\]
and the assumed lower bound on \(n\) ensures \(t_n\le1/2\). On the event
\(|D_n-1|\le t_n\le1/2\),
\[
  \|\hatA_\psi-A^\mu_\psi\|_{\mathrm{op}}
  =
  \left\|\frac{\tildeA_\psi}{D_n}-A^\mu_\psi\right\|_{\mathrm{op}}
  \le
  2\|\tildeA_\psi-A^\mu_\psi\|_{\mathrm{op}}
  +
  2|D_n-1|\,\|A^\mu_\psi\|_{\mathrm{op}}.
\]
The same argument gives
\[
  \|\hatb_\psi-b^\mu_\psi\|_2
  \le
  2\|\tildeb_\psi-b^\mu_\psi\|_2
  +
  2|D_n-1|\,\|b^\mu_\psi\|_2.
\]
Finally, \(\|A^\mu_\psi\|_{\mathrm{op}}\le K^2(1+\gamma)\) and
\(\|b^\mu_\psi\|_2\le KR_{\max}\). Absorbing constants into a universal
constant \(C\) proves the stated inequalities.
\end{proof}

\begin{proof}[Proof of Corollary~\ref{cor:tree}]
For one-hot tree features, the normal equations are
\[
  D_\mu(I-\gamma P^\mu)q=D_\mu r^\mu,
\]
which reduces to \(q=r^\mu+\gamma P^\mu q\) on leaves with positive
\(\mu\)-mass. The empirical equation is the same with
\((\hatr^\mu,\hatP^\mu)\), where
\[
  \hatr^\mu_\ell
  =
  \frac{n^{-1}\sum_i w_i\one\{\phi_{\cT}(X_i)=\ell\}R_i}
       {\hat\mu_\ell},
  \qquad
  \hatP^\mu_{\ell m}
  =
  \frac{n^{-1}\sum_i w_i\one\{\phi_{\cT}(X_i)=\ell\}
    p_{\cT,m}(S'_i)}
       {\hat\mu_\ell}.
\]
For illustration, let
\[
  N^r_\ell=\E_\nu[w\one\{\phi_{\cT}(X)=\ell\}R],
  \qquad
  \hat N^r_\ell=n^{-1}\sum_iw_i\one\{\phi_{\cT}(X_i)=\ell\}R_i.
\]
On \(\hat\mu_\ell\ge\mu_\ell/2\),
\[
  |\hatr^\mu_\ell-r^\mu_\ell|
  =
  \left|\frac{\hat N^r_\ell}{\hat\mu_\ell}
  -\frac{N^r_\ell}{\mu_\ell}\right|
  \le
  \frac{2|\hat N^r_\ell-N^r_\ell|}{\mu_\ell}
  +
  \frac{2R_{\max}|\hat\mu_\ell-\mu_\ell|}{\mu_\ell}.
\]
The same ratio calculation applies entrywise to \(\hatP^\mu_{\ell m}\), with
\(R_{\max}\) replaced by \(1\). Thus the displayed assumptions
\(\linf{\hatr^\mu-r^\mu}\le\varepsilon_r\) and
\(\|\hatP^\mu-P^\mu\|_\infty\le\varepsilon_P\) are the corresponding
self-normalized leaf-moment errors.

Subtracting the two Bellman equations and using
\(\|(I-\gamma\hatP^\mu)^{-1}\|_\infty\le(1-\gamma)^{-1}\) and
\(\linf{q}\le R_{\max}/(1-\gamma)\) gives the displayed bound. The leaf-mass
condition follows because
\[
  \Pr\{|\hat\mu_\ell-\mu_\ell|>\mu_{\min}/2\}
  \le
  2\exp\{-n\mu_{\min}^2/(2W^2)\}
\]
for each leaf by Hoeffding, and a union bound over \(L\) leaves gives the
claimed sample-size scaling.
\end{proof}

\begin{proof}[Proof of Proposition~\ref{prop:weight_perturbation}]
Apply the same perturbation argument as in
Theorem~\ref{thm:feature_perturbation} with
\((A,b)=(A^w,b^w)\) and \((\hatA,\hatb)=(A^{\hat w},b^{\hat w})\).
\end{proof}

\begin{proof}[Proof of Proposition~\ref{prop:validation}]
Condition on the candidate set and on all non-validation data used to compute
the provisional residual functions. For a fixed candidate,
\(w_i\delta_{\psi,i}^2\in[0,WB_\delta^2]\).
Hoeffding's inequality and a union bound over \(\mathfrak P\) imply a uniform
deviation bound of
\[
  WB_\delta^2
  \sqrt{\frac{\log(2|\mathfrak P|/\delta)}{2n_{\mathrm{val}}}}.
\]
Call this deviation \(\varepsilon_{\mathrm{val}}\). Let
\(\psi^\star\) minimize the population penalized risk over \(\mathfrak P\).
On the uniform event,
\[
  \begin{aligned}
  L_\mu(\hat\psi)+\lambda\dim(\spann\hat\psi)
  &\le
  \widehat L_{\mathrm{val}}(\hat\psi)
  +\lambda\dim(\spann\hat\psi)
  +\varepsilon_{\mathrm{val}}\\
  &\le
  \widehat L_{\mathrm{val}}(\psi^\star)
  +\lambda\dim(\spann\psi^\star)
  +\varepsilon_{\mathrm{val}}\\
  &\le
  L_\mu(\psi^\star)
  +\lambda\dim(\spann\psi^\star)
  +2\varepsilon_{\mathrm{val}}.
  \end{aligned}
\]
This proves the result. If the empirical score is divided by the validation
denominator \(D_{\mathrm{val}}\), the same argument applies on the event
\(|D_{\mathrm{val}}-1|\le 1/2\), with the displayed deviation multiplied by a
universal constant.
\end{proof}

\section{Implementation Details}
\label{app:implementation}

\paragraph{Candidate splits.}
For continuous variables, candidate thresholds are empirical quantiles inside
each leaf after enforcing minimum raw count, weighted mass, and ESS. For
categorical variables and finite actions, the implementation considers binary
subsets when arity is small and one-vs-rest splits otherwise. The action is
treated as part of the state-action feature vector.

\paragraph{Rank reduction.}
For the theorem-clean estimator, forest features are reduced on an independent
rank fold by QR with column pivoting or SVD using a fixed tolerance. The
retained feature map is then held fixed for final estimation. A practical
implementation may run the same QR/SVD on the estimation design because it uses
no rewards, but this should be reported as a data-dependent numerical
preprocessing step. A ridge fallback is allowed only as a diagnostic because it
changes the projected equation.

\paragraph{Residual scoring.}
The residual score in \eqref{eq:score} uses held-out weighted transitions. A
candidate's provisional value is estimated on the growth-fit fold, then
residuals are computed on growth-score or validation transitions by assigning
current and next state-action pairs to candidate features and weighting by
\(\hat w_i\).

\section{Experiment Templates}
\label{app:experiment_templates}

\paragraph{Diagnostic MDP template.}
Use a two-dimensional state \(S=(S_1,S_2)\in[-1,1]^2\) and two actions. Let one
coordinate primarily determine immediate rewards and the other determine which
high- or low-value region the target policy reaches next. Display true
\(Q^\pi\), CART/RF supervised partitions or weights, Tree-FQE/RF-FQE behavior,
BAT partitions, and BAF forest weights at matched feature budgets.

\paragraph{Navigation template.}
Use a grid-auditable continuous navigation task with five actions. Compute
reference \(Q^\pi\) by dynamic programming on a fine grid. Draw offline
transitions from behavior policies with increasing mismatch from the target
stationary distribution. Report errors under both the target and behavior
distributions to show when ordinary supervised criteria optimize the wrong
geometry.

\end{document}
